\section{Cotangent resolvents: a nonpolynomial Stieltjes generator}
\label{res:section}

Question Q8 of the incoming \emph{Exact Resonant Identities} report
\cite{ghp:ExactResonant} asks for a nonpolynomial extension of its
cotangent moment calculus, including the local subtraction before
parameter expansion. We give such an extension for every rational
function of the cotangent bounded at infinity and with no real pole. A single
resolvent supplies the answer; differentiation gives all repeated
poles, and partial fractions give the general rational case. The
answer involves ordinary polylogarithms and their order derivatives.
The Fourier identity used in the proof is classical. The contribution
here is the completed rational-kernel identity, its precise endpoint
correction at every Stieltjes and argument-derivative order, and its
compatibility with continuous polygamma orders and primitives.

\subsection{Coordinates, branches, and finite parts}

Write
\begin{equation}\label{res:kernel}
 K(x)=\pi\cot(\pi x),\qquad
 R_z(x)=\frac1{K(x)-z},\qquad 0<x<1,\quad z\in\mathbb C\setminus\mathbb R.
\end{equation}
Although $K$ has endpoint poles, $R_z$ extends to a smooth periodic
function, with $R_z(0)=R_z(1)=0$. On either open half-plane put
\begin{equation}\label{res:coordinates}
 \epsilon=\operatorname{sgn}(\Im z),\quad
 a=z+\epsilon\pi i,\quad
 q=\frac{z-\epsilon\pi i}{z+\epsilon\pi i},\quad
 c=\frac{2\epsilon\pi i}{a^2},\quad
 \sigma=-2\epsilon\pi i.
\end{equation}
Thus $|q|<1$, and throughout the section
\[
 \Log\sigma=\log(2\pi)-\epsilon i\pi/2.
\]
The letter $a$ in this section is only the denominator in
\eqref{res:coordinates}; it is not a Hurwitz endpoint.
To remove an otherwise unnecessary exceptional case, define
\begin{equation}\label{res:normalizedpolylog}
 \Phi_*(s,q)=\frac{\Li_s(q)}q
       =\sum_{n\ge1}\frac{q^{n-1}}{n^s},\qquad
 \Phi_*(s,0)=1.
\end{equation}
This function is entire in $s$ and holomorphic for $|q|<1$.
Its order derivatives at $q=0$ are zero except at order zero.

We use the conventional generalized Stieltjes functions
\begin{equation}\label{res:stieltjes}
 \zeta(1-s,x)=-\frac1s+
       \sum_{m\ge0}\gamma_m(x)\frac{s^m}{m!},\qquad
 \gamma_0(x)=-\psi(x).
\end{equation}
A superscript $(p)$ on $\gamma_m$ or $\psi$ means an argument
derivative. A superscript $[j]$ on $\Li_s$ or $\Phi_*$ means an order
derivative. These two operations must be distinguished.

Every finite part below is taken in the original variable $x$:
$\operatorname{FP}_x\int_0^1 f(x)\,dx$ is the coefficient of
$\varepsilon^0(\log\varepsilon)^0$ in
$\int_\varepsilon^{1-\varepsilon}f(x)\,dx$. Equivalently, subtract
the finitely many nonintegrable endpoint terms, integrate the
remainder, and add the constant terms of the explicit subtractions.
For the resolvent integrals only the endpoint $x=0$ can be singular.
The Stieltjes shift formula gives
\begin{equation}\label{res:singularpart}
 \gamma_m(x)=\frac{\log^m x}{x}+\gamma_m(x+1).
\end{equation}
Consequently all the finite parts used here exist. In particular,
\[
 M_{m,p}(z):=\operatorname{FP}_x\int_0^1\gamma_m^{(p)}(x)R_z(x)\,dx
\]
is an ordinary convergent integral when $p=0$.

\subsection{The complete all-index identity}

Set
\begin{equation}\label{res:BE}
 B_\epsilon(s)=\Gamma(1+s)\sigma^{-s},\qquad
 E_p(s)=\prod_{j=1}^p\left(1-\frac{s}{j}\right),\qquad E_0(s)=1.
\end{equation}

\begin{theorem}[Completed cotangent--Stieltjes resolvent]
\label{res:main}
For all integers $m,p\ge0$ and $z\notin\mathbb R$,
\begin{equation}\label{res:master}
 \boxed{\displaystyle
 M_{m,p}(z)=m!c\sigma^p[s^{m+1}]
 \left\{B_\epsilon(s)\Phi_*(s-p,q)
                   -E_p(s)\Phi_*(-p,q)\right\}.}
\end{equation}
The expression in braces vanishes at $s=0$. The identity is
holomorphic in $z$ on each half-plane, including $z=\epsilon i\pi$.
For $p\ge1$ the polynomial $E_p$ is the full local endpoint
correction. Its contribution vanishes whenever $m\ge p$.
\end{theorem}

The identity can be expanded without any coefficient notation.
Let $g_j=B_\epsilon^{(j)}(0)$ and $e_{p,j}=E_p^{(j)}(0)$. Then
\begin{equation}\label{res:expanded}
 M_{m,p}(z)=\frac{c\sigma^p}{m+1}
 \left\{\sum_{j=0}^{m+1}\binom{m+1}{j}
     g_{m+1-j}\Phi_*^{[j]}(-p,q)
                 -e_{p,m+1}\Phi_*(-p,q)\right\}.
\end{equation}
The coefficients $g_j$ are complete exponential Bell polynomials
with logarithmic derivatives
\[
 b_1=-\gamma-\Log\sigma,\qquad
 b_j=(-1)^j(j-1)!\zeta(j)\quad(j\ge2),
\]
so $g_0=1$, $g_1=b_1$, $g_2=b_1^2+\zeta(2)$, and
$g_3=b_1^3+3b_1\zeta(2)-2\zeta(3)$.
Likewise
\begin{equation}\label{res:contactcoeff}
 e_{p,j}=(-1)^j j!\,
 e_j\left(1,\frac12,\ldots,\frac1p\right),
\end{equation}
where $e_j$ on the right is the elementary symmetric polynomial;
it is zero for $j>p$. Thus the contact terms use ordinary harmonic
numbers and their finite Bell-polynomial combinations.

\begin{proof}[Proof of Theorem~\ref{res:main}]
It is enough to prove the upper-half-plane case; reversing the
Fourier signs gives the other case with exactly the branches in
\eqref{res:coordinates}. Put $w=e^{2\pi ix}$. Direct algebra gives
the uniformly convergent Fourier expansion
\begin{equation}\label{res:fourierkernel}
 R_z(x)=-\frac1a+c\sum_{n\ge1}q^{n-1}e^{2\pi inx}.
\end{equation}
For example, at $z=i\pi$ this reduces to
$(e^{2\pi ix}-1)/(2\pi i)$. The classical Hurwitz Fourier formula
\cite[\S25.11]{DLMF} implies, initially in a sufficiently large
right half-plane,
\[
 \int_0^1\zeta(1-s,x)\,dx=0,\qquad
 \int_0^1\zeta(1-s,x)e^{2\pi inx}\,dx
       =\Gamma(s)\sigma^{-s}n^{-s}.
\]
All termwise integrations are legitimate there because of the
geometric factor $q^n$. Argument differentiation multiplies the
negative Fourier mode by $\sigma n$. For $\Re s>p+1$, this gives
\begin{equation}\label{res:spectralintegral}
 J_p(s,z):=\int_0^1\partial_x^p\zeta(1-s,x)R_z(x)\,dx
       =c\Gamma(s)\sigma^{p-s}\Phi_*(s-p,q).
\end{equation}
This equation is then an identity of meromorphic continuations.

We now relate that continuation to the pointwise finite part;
this is the step that cannot be replaced by substitution at $s=0$.
Write $r_j(z)=[x^j]R_z(x)$. For $p\ge1$, the sole local term
whose moving exponent crosses $-1$ at $s=0$ is
\begin{equation}\label{res:localmoving}
 r_p(z)A_p(s)x^{s-1},\qquad
 A_p(s)=\prod_{j=1}^p(s-j).
\end{equation}
Indeed the singular part of
$\partial_x^p\zeta(1-s,x)$ is
$A_p(s)x^{s-p-1}$. The analytic part from $\zeta(1-s,x+1)$
has no pole in $s$ after an argument derivative. Every other
Taylor term of $R_z$ has an exponent separated from $-1$ near
zero. For such a term, finite-part integration commutes with
coefficient extraction. The term \eqref{res:localmoving} is different:
its analytically continued integral is $r_pA_p(s)/s$, whereas the
finite part of every coefficient $x^{-1}\log^j x$ on $(0,1)$ is zero.
It follows, with the entire amplitude retained, that
\begin{equation}\label{res:completion}
 M_{m,p}(z)=m![s^m]
            \left(J_p(s,z)-\frac{r_p(z)A_p(s)}s\right),\quad p\ge1.
\end{equation}
One can make this argument on $(0,\delta)$ with a finite Taylor
subtraction and a normally integrable remainder. The compensating
integrals on $(\delta,1)$ cancel the artificial splitting point.
This also proves all holomorphy and coefficient interchanges used here.

Differentiating \eqref{res:fourierkernel} at $x=0$ gives
\begin{equation}\label{res:rj}
 r_p(z)=\frac{c(-\sigma)^p}{p!}\Phi_*(-p,q),\qquad p\ge1.
\end{equation}
Since $A_p(s)=(-1)^pp!E_p(s)$, substitution in
\eqref{res:completion} proves \eqref{res:master} for $p\ge1$.

For $p=0$, the product $\gamma_m R_z$ is integrable by
\eqref{res:singularpart}. The pole in \eqref{res:stieltjes} is
independent of $x$, so
\[
 \sum_{m\ge0}M_{m,0}(z)\frac{s^m}{m!}
       =J_0(s,z)+\frac1s\int_0^1R_z(x)\,dx
       =J_0(s,z)-\frac1{as}.
\]
But $c\Phi_*(0,q)=1/a$. This is exactly
\eqref{res:master} with $E_0=1$. Thus the same formula accounts
for both the global Hurwitz pole at $p=0$ and the local moving
endpoint singularity at $p\ge1$.
\end{proof}

\subsection{Explicit examples and a check of the correction}

At the first Stieltjes index the formula is
\begin{equation}\label{res:polygamma}
 \operatorname{FP}_x\int_0^1\psi^{(p)}(x)R_z(x)\,dx
  =c\sigma^p\left\{(\gamma+\Log\sigma-H_p)\Phi_*(-p,q)
                                      -\Phi_*^{[1]}(-p,q)\right\},
\end{equation}
with $H_0=0$. The harmonic number is essential whenever $p>0$.
For $m=1$ the complete formula is
\begin{align}\label{res:gamma1}
 M_{1,p}(z)=\frac{c\sigma^p}{2}\bigl\{
  &\Phi_*^{[2]}(-p,q)-2(\gamma+\Log\sigma)\Phi_*^{[1]}(-p,q)\notag\\
  &+\bigl((\gamma+\Log\sigma)^2+\zeta(2)
                      -H_p^2+H_p^{(2)}\bigr)\Phi_*(-p,q)\bigr\}.
\end{align}
Here $H_0^{(2)}=0$. The correction in this equation vanishes at
$p=0,1$, as the general threshold predicts.

At $z=i\pi$, the normalized polylogarithm is identically one,
and all positive order derivatives disappear. In particular,
\begin{equation}\label{res:explicitcheck}
 \operatorname{FP}_x\int_0^1
       \frac{\psi'(x)}{\pi\cot(\pi x)-i\pi}\,dx
       =1-\gamma-\log(2\pi)+\frac{i\pi}{2}.
\end{equation}
There is a direct independent verification of its constant term.
At this $z$, $R'_z=e^{2\pi ix}$ and the finite boundary value of
$[\psi R_z]_0^1$ is $1$, because $\psi(x)R_z(x)\to-1$ as
$x\downarrow0$. Integration by parts therefore supplies precisely
the $1$ in \eqref{res:explicitcheck}. Omitting the local correction
would miss it.

The local coefficients themselves require no special functions:
\begin{equation}\label{res:riccati}
 R_z'=1+2zR_z+(z^2+\pi^2)R_z^2,\qquad R_z(0)=0.
\end{equation}
Thus $r_1=1$, $r_2=z$, $r_3=z^2+\pi^2/3$, and
$r_4=z^3+2\pi^2z/3$. In general $r_p$ is a polynomial in $z$
of degree $p-1$. The difference between the raw spectral coefficient
and the stated finite part is therefore a finite polynomial in $z$,
not an additional unevaluated integral.

For a real parameter $b>0$, taking the real and imaginary parts of
the $p=m=0$ formula at $z=ib$ gives ordinary real integrals. For example,
with $q=(b-\pi)/(b+\pi)$,
\begin{align}
 \int_0^1\frac{\gamma_0(x)K(x)}{K(x)^2+b^2}\,dx
       &=\frac{\pi}{2(b+\pi)},\label{res:realexample1}\\
 \int_0^1\frac{\gamma_0(x)}{K(x)^2+b^2}\,dx
       &=\frac{\gamma+\log(2\pi)
                   -(1-q)\Phi_*^{[1]}(0,q)}{b(b+\pi)}.
       \label{res:realexample2}
\end{align}
At $b=3\pi$, the last numerator becomes
$\gamma+\log(2\pi)-\Li_0^{[1]}(1/2)$.
Equation \eqref{res:realexample1} also follows directly from the
digamma reflection formula, providing a separate check of the phase.

\subsection{Repeated poles and the complete bounded-rational family}

\begin{corollary}[All repeated resolvents]\label{res:repeated}
For every integer $r\ge1$,
\begin{equation}\label{res:repeatedformula}
 \operatorname{FP}_x\int_0^1
       \frac{\gamma_m^{(p)}(x)}{(K(x)-z)^r}\,dx
       =\frac1{(r-1)!}\partial_z^{r-1}M_{m,p}(z).
\end{equation}
This is an ordinary absolutely convergent integral if $r>p$.
Every rational function $Q(u)$ bounded at infinity and with no real pole admits
a finite reduction of
$\operatorname{FP}_x\int_0^1\gamma_m^{(p)}(x)Q(K(x))\,dx$
to the right-hand sides of \eqref{res:repeatedformula}.
\end{corollary}

\begin{proof}
Parameter differentiation commutes with the finitely many local
subtractions on compact subsets away from the real poles. Since
$\partial_z^{r-1}R_z=(r-1)!R_z^r$, the formula follows.
Near zero the integrand has size
$O(x^{r-p-1}(1+|\log x|^m))$, proving the convergence assertion.
Equivalently, the resonant coefficient $[x^p]R_z^r$ is zero
when $r>p$, so all contact terms disappear from the differentiated
formula.

Partial fractions express $Q$ as a constant plus finitely many
repeated resolvents. The constant contributes zero:
\[
 \operatorname{FP}_x\int_0^1\gamma_m^{(p)}(x)\,dx=0.
\]
For $p=0$ this follows by integrating the Hurwitz spectral family
and matching its constant and endpoint pole. For $p\ge1$ it
follows from the finite part of the endpoint difference of
$\gamma_m^{(p-1)}$, using \eqref{res:singularpart}. The analytic
parts at $0$ and $1$ agree, and the singular part has zero constant.
This proves the complete finite reduction.
\end{proof}

Repeated $z$ derivatives introduce no new transcendental function
class: $\partial_q\Li_s(q)=\Li_{s-1}(q)/q$, and $q(z)$ is rational.
The reduction therefore stays within finitely many polylogarithmic
order derivatives at explicit arguments. This proves the rational
sector of the proposed nonpolynomial extension. It does not assert
the same finite closure for arbitrary complex powers of $K-z$.

\section{Continuous polygamma orders, Gamma primitives, and boundary values}
\label{res:primitive-section}

\subsection{An entire transform and its integer-order anomaly}

We use the established generalized polygamma function of Espinosa
and Moll \cite{EM}, denoted here by $\Psi(v,x)$ to distinguish it
from the ordinary digamma:
\begin{equation}\label{res:EM}
 \Psi(v,x)=e^{-\gamma v}\partial_v
       \left(e^{\gamma v}\frac{\zeta(1+v,x)}{\Gamma(-v)}\right).
\end{equation}
The removable values in this definition are understood. It is entire
in $v$, satisfies $\partial_x\Psi(v,x)=\Psi(v+1,x)$, and equals
$\psi^{(p)}(x)$ when $v=p\ge0$. These are properties of the
classical generalized function, not new definitions of polygamma.

\Needspace{18\baselineskip}
\begin{theorem}[Continuous-order resolvent and finite-part conversion]
\label{res:continuous}
For $\Re v<1$ and $z\notin\mathbb R$, the following ordinary
absolutely convergent integral holds:
\begin{equation}\label{res:continuousformula}
 \mathcal C(v,z):=\int_0^1\Psi(v,x)R_z(x)\,dx
   =c\sigma^v\left\{(\gamma+\Log\sigma)\Phi_*(-v,q)
                                         -\Phi_*^{[1]}(-v,q)\right\}.
\end{equation}
The right side is entire in $v$. At a nonnegative integer $p$,
its continuation and the pointwise finite part are related by
\begin{equation}\label{res:anomaly}
 \operatorname{FP}_x\int_0^1\psi^{(p)}(x)R_z(x)\,dx
    =\mathcal C(p,z)-(-1)^pp!H_p r_p(z),
\end{equation}
where the correction is zero for $p=0$ without defining $r_0$.
\end{theorem}

\begin{proof}
The singular part of \eqref{res:EM} is a linear combination of
$x^{-v-1}$ and $x^{-v-1}\log x$ with entire coefficient functions
of $v$. Multiplication by $R_z=O(x)$ gives locally uniform
integrability when $\Re v<1$. Thus the integral is holomorphic
in that half-plane, including the removable point $v=0$.

For $\Re v<0$, divide \eqref{res:spectralintegral} with $p=0$
by $\Gamma(s)$ and set $s=-v$. Applying
$e^{-\gamma v}\partial_v e^{\gamma v}$ yields
\eqref{res:continuousformula}. Holomorphic continuation gives
the entire stated ordinary-integral strip. Finally, compare the
right side at $v=p$ with \eqref{res:polygamma}, and use
$c\sigma^p\Phi_*(-p,q)=(-1)^pp!r_p(z)$.
This proves the exact conversion formula.
\end{proof}

The last statement is a concrete instance of noncommutation between
specializing a parameter and taking an endpoint finite part. It is
not a defect of the generalized polygamma function. At $p=1$ the
correction in \eqref{res:anomaly} is $+1$ for every nonreal $z$.
At $p=2$ it is $-3z$. These exact corrections are useful checks
when a spectral formula is integrated into a symbolic system.

Integration by parts gives a compatible differential identity for
the entire transform:
\begin{equation}\label{res:continuous-recurrence}
 \mathcal C(v+1,z)
      =-2z\mathcal C(v,z)-(z^2+\pi^2)\partial_z\mathcal C(v,z).
\end{equation}
Initially take $\Re v<0$, so the boundary term vanishes and
$\int_0^1\Psi(v,x)\,dx=0$, and use \eqref{res:riccati}.
Both sides are entire in $v$, which proves the identity everywhere.
For the pointwise finite-part moments
$F_p(z)=\operatorname{FP}_x\int\psi^{(p)}R_z$, the corresponding
identity instead reads
\begin{equation}\label{res:pointwise-recurrence}
 F_{p+1}(z)=-2zF_p(z)-(z^2+\pi^2)F_p'(z)
                         +(-1)^pp!r_{p+1}(z).
\end{equation}
The last term is exactly the finite part of $[\psi^{(p)}R_z]_0^1$.
Thus the derivative hierarchy, the endpoint subtraction, and the
polylogarithmic formula agree at every order.

\subsection{Every balanced Gamma primitive}

For $r\ge0$, let
\begin{equation}\label{res:balanced}
 \mathcal B_r(x)=\frac{\zeta'(-r,x)+H_r\zeta(-r,x)}{r!}.
\end{equation}
These are the standard balanced negapolygamma functions
$\Psi(-r-1,x)$ \cite{EM}. Direct differentiation gives
\[
 \mathcal B_0(x)=\log\Gamma(x)-\tfrac12\log(2\pi),\qquad
 \mathcal B_r'(x)=\mathcal B_{r-1}(x)\quad(r\ge1).
\]
The additive constants have therefore been specified completely.

\begin{corollary}[The rational-kernel primitive ladder]
\label{res:balancedcor}
For every $r\ge0$,
\begin{equation}\label{res:balancedmoment}
 \int_0^1\mathcal B_r(x)R_z(x)\,dx
    =c\sigma^{-r-1}
       \left\{(\gamma+\Log\sigma)\Phi_*(r+1,q)
                                      -\Phi_*^{[1]}(r+1,q)\right\}.
\end{equation}
All these integrals are ordinary convergent integrals.
\end{corollary}

This follows from \eqref{res:continuousformula} at $v=-r-1$.
There is also a useful direct check. Differentiate
\eqref{res:spectralintegral} in $s$ at $s=r+1$. The derivative of
$\Gamma(s)$ supplies $H_r-\gamma$; the $H_r\zeta(-r,x)$ term
in \eqref{res:balanced} cancels $H_r$ exactly. That cancellation
explains why the same coefficient $\gamma+\Log\sigma$ works at
every primitive level.

In particular,
\begin{equation}\label{res:loggammamoment}
 \int_0^1\log\Gamma(x)R_z(x)\,dx
 =-\frac{\log(2\pi)}{2a}
   +\frac{\Phi_*^{[1]}(1,q)
                 -(\gamma+\Log\sigma)\Phi_*(1,q)}{a^2}.
\end{equation}
At the next level, the standard Barnes identity
\cite{Vardi}
\[
 \zeta'(-1,x)=\zeta'(-1)-\log G(x)+(x-1)\log\Gamma(x)
\]
shows that \eqref{res:balancedmoment} is an explicit formula for
the corresponding normalized Barnes-$G$ combination, since
$\zeta(-1,x)=-B_2(x)/2$. Here $G$ is the standard Barnes function, with its usual asymptotic
normalization, $G(x+1)=\Gamma(x)G(x)$, and $G(1)=1$. The
functional equation and the value at $1$ alone would not exclude
a multiplicative periodic factor. This supplies an exact primitive
family relevant to Q9 in the source, without claiming to evaluate
every unbalanced product containing $G$.

\subsection{Boundary values and recovery of the Stieltjes functions}

The resolvent also permits a useful inversion check. Given
$u\in\mathbb R$, let $x_u\in(0,1)$ be the unique point with
$K(x_u)=u$. For $m\ge0$ the ordinary moment $M_{m,0}$ has
one-sided boundary values satisfying
\begin{equation}\label{res:jump}
 M_{m,0}(u+i0)-M_{m,0}(u-i0)
       =\frac{2\pi i}{u^2+\pi^2}\gamma_m(x_u).
\end{equation}
Indeed the substitution $y=K(x)$ gives the Cauchy integral
\[
 M_{m,0}(z)=\int_{-\infty}^{\infty}
      \frac{\gamma_m(x_y)}{(y^2+\pi^2)(y-z)}\,dy.
\]
Its tail is integrable after the factor $(y-z)^{-1}$ is included,
and its density is smooth locally at $y=u$. Splitting off that
local value and integrating $1/(y-z)$ proves
\eqref{res:jump}; this is the elementary Cauchy jump formula.
The average of the two boundary values gives the corresponding
principal-value integral at the interior pole.

On the boundary, $q=e^{-2\pi ix_u}$ in the upper half-plane and
$q=e^{2\pi ix_u}$ in the lower one. Thus the two versions of
\eqref{res:master} recover the classical Hurwitz--polylogarithm
connection and all its Stieltjes order derivatives. This boundary
interpretation is a consistency check and an exact inversion formula;
it is not claimed as a new Hurwitz functional equation.

\subsection{Verification and the precise new scope}

The program \src{code/verify_resolvent.py} checks the Fourier and
Riccati expressions for the local coefficients by exact rational
function arithmetic, including both the Gamma polynomial and its
harmonic correction. Its numerical checks compare the formulas with
direct quadrature of digamma and generalized Stieltjes functions,
local-subtraction quadrature for the finite parts, convergent
integrals involving $\gamma_1'$, and the first three balanced
Gamma primitives. Complex-order tests use the defining Hurwitz
formula \eqref{res:EM}. The numerical and exact records are distinct.

The all-index rational closure answers the specified rational sector
of the source's nonpolynomial question. The local correction
mechanism extends the repository's established finite-part calculus;
it is not presented as an independently invented regularization.
The continuous-order conversion and the primitive ladder explain
which explicit boundary terms must accompany the resulting identities.
